\section{Threats to Validity}
\label{sec:threats}

\paragraph{Internal threats.} These threaten whether \tool's own results
are correct.
\begin{itemize}
  \item \emph{Tuning on seen data.} The deref-after-release suppression and
    the \texttt{NULL}-guard pruning (Section~\ref{sec:step2}) were both
    added after reading real output on OpenSSL. The second was added
    specifically after it produced one infeasible-path false positive on
    the master scan. The master-snapshot numbers in
    Table~\ref{tab:rq1-wholetree} are therefore not a blind test of those
    two refinements, though the rule itself (Section~\ref{sec:rule}) was
    not changed after seeing the KDF result in
    Observation~\ref{obs:rq1-conformance}.
  \item \emph{Release detection is name-based.} A release site is any call
    matching \texttt{(?i)free}. This cannot see reference counts: one site
    in an earlier snapshot, \texttt{load\_common}, frees only a structural
    \texttt{ENGINE} reference and is a confirmed false positive under
    manual review.
  \item \emph{Parser coverage.} About 39\% of scanned files have a
    \texttt{tree-sitter} parse error somewhere, almost always from an
    unexpanded macro. Only 7 to 8 release sites per snapshot sit inside a
    function with a parse error, and no flagged violation in
    Section~\ref{sec:eval} is among them. A site inside an unparseable
    function is invisible to \tool, though, and would not appear in any
    count we report.
  \item \emph{Our own triage.} The labels behind
    Observation~\ref{obs:rq1-mastertriage} (2 potential recurrence, 8
    convention deviation) were assigned by the people who built
    \tool. We had one independent party re-derive the reachability
    argument for the 2 potential-recurrence sites from scratch against
    real source. (The party was a generative AI assistant, not involved
    in building the rule or running the original scan; see
    \texttt{results/reachability\_rq1\_candidates.md}.) It found and
    corrected a factual error in our original note (Section~\ref{sec:threats}
    construct threats, below). The same party also re-checked the 2
    sites we had labelled "unclear" pending a caller read
    (\texttt{results/unclear\_sites\_recheck.md}). Both reclassify as
    convention deviation once the one real in-tree caller
    (\texttt{apps/cmp.c}) is read, since it discards the dangling output
    on every failure path rather than acting on it; the failure
    branch itself needs the same unreachable-without-a-race precondition
    as the two potential-recurrence sites. Every one of the 10
    whole-tree flags now has an independently re-derived reason, not
    just our own. That one re-check already changed a conclusion, which
    is itself evidence for, not against, this threat: a human second
    labeller without knowledge of the rule's derivation remains undone
    and could still disagree with any of these 10 labels.
\end{itemize}
What remains: a human independent second labeller on a sample of the
triaged sites, and a parser with macro expansion, are both future work.

\paragraph{External threats.} These threaten whether the results
generalize past our six snapshots and three CVEs.
\begin{itemize}
  \item \emph{One rule family generalizes within itself, not across
    invariant shapes.} Protocol B, derived from CVE-2022-4450
    (Observation~\ref{obs:rq1-conformance},
    Observation~\ref{obs:rq3-timetravel}), also reproduces an unrelated
    fix and flags a site a later commit fixed. That is two data points,
    though, not a sample large enough for a detection rate. Protocol A, a
    differently-shaped rule for CVE-2023-0215 (Section~\ref{sec:rq5}),
    distinguishes its own motivating fix, but its whole-tree run
    (Observation~\ref{obs:rq5-protocola-master}) is 0 for 11. Every
    flagged site is a false positive, because the rule has no notion of
    which released object is externally owned. That is a real negative
    result, not a bug to patch away: a single-object release/detach rule
    does not, by itself, generalize the way Protocol B's did.
  \item \emph{Non-independent snapshots.} The six snapshots in
    Table~\ref{tab:rq1-wholetree} share most of their long-lived
    functions across revisions, so the effective sample is closer to 20
    distinct release sites than to 20~$\times$~6.
  \item \emph{One project.} All subjects (Section~\ref{sec:subjects}) are
    OpenSSL. We make no claim about another C codebase with different
    allocation and ownership conventions.
  \item \emph{Scope of what is tracked.} \tool only examines parameters of
    pointer-to-pointer type or deeper. An out-parameter hidden behind a
    struct member, such as \texttt{ctx->buf}, is invisible to it, however
    it is released or reset.
\end{itemize}
What remains: more historical pairs, a second project, tracking
struct-member out-parameters, and an ownership/escape analysis that could
make Protocol A's whole-tree run precise, are future work.

\paragraph{Construct threats.} These threaten whether what we measure is
what we claim.
\begin{itemize}
  \item \emph{A violation is a source-pattern property, not a proven
    vulnerability.} \tool reports that some path reaches function exit
    with \texttt{*P} still released; it does not establish that the path
    is reachable at runtime. Of the 2 potential-recurrence sites in
    Observation~\ref{obs:rq1-mastertriage}, both need the same precondition
    once independently re-derived against real source
    (\texttt{results/reachability\_rq1\_candidates.md}): the tracked
    object's state must change between two reads a single-threaded call
    proves equal, which only an unsynchronized second thread racing the
    same object can cause. No in-tree caller creates this precondition. We
    also found and corrected one factual error in our own earlier note
    (a claimed lock/atomic-op failure that is dead code on every platform
    this codebase builds for). We report both as candidates needing a
    multithreaded fault-injection test, not as bugs.
  \item \emph{A clean run on a general tool means absence of a diagnostic,
    not absence of a bug.} Table~\ref{tab:rq2-comparison}'s \xmark{} cells
    mean the tool emitted nothing for that translation unit under that
    configuration (Section~\ref{sec:baselines}). Clang Static Analyzer, in
    particular, analyzed only the one affected translation unit per
    revision, not the whole build. A clean run is therefore a property of
    that unit and configuration, not a property of the tool in general.
  \item \emph{RQ2's baseline is not \tool.} Table~\ref{tab:rq2-comparison}
    compares Clang, Coccinelle, and Cppcheck against the Part 1 pattern
    matcher, not against Protocol B. A reader who conflates the two would
    credit Protocol B with a result it did not produce; we repeat this
    distinction at the point it matters (Section~\ref{sec:rq2}) to avoid
    that.
\end{itemize}
